\section{Conclusion}
In this work, we address two persistent bottlenecks in visuo-tactile manipulation: the scarcity of large-scale, task-diverse training data and the lack of methods that model contact dynamics explicitly while supporting closed-loop tactile control.
On the data side, OmniViTac comprises $21{,}000+$ temporally aligned trajectories spanning $86$ contact-rich tasks organized into six physics-grounded interaction patterns. Systematic analysis of the dataset identifies two structural properties of tactile signals in manipulation, spatial locality, and contact-driven dynamics, which guided the architectural design of OmniVTA.
On the methodology side, OmniVTA integrates self-supervised implicit tactile representation learning, short-horizon visuo-tactile world modeling, contact-aware multimodal fusion, and high-frequency reflexive control within a unified framework. Real-robot experiments across all six interaction categories demonstrate that OmniVTA consistently outperforms existing baselines, with clear advantages in generalization to unseen objects and geometric configurations and robustness to external disturbances.

\section{Acknowledgments}
We want to thank Xiangyu Chen, Yating Feng, Zhijie Yan, Xiang Xiao, Jingxiang Guo, and Yuxing Qin for their efforts in building the OmniViTac dataset.
% \subsection{Limitation and future work}
% This work has several limitations. While OmniViTac provides a comprehensive benchmark for single-arm, gripper-based tactile manipulation, it does not currently encompass the dual-arm setting or other end-effector types, such as dexterous hands. Future work will investigate scaling the world model with larger and more diverse data, extending the framework to dexterous hands and dual-arm manipulation, and exploring cross-embodiment transfer of the visuo-tactile representations.